\section{Constructing the mixed automorphic seed}
\label{sec:symplectic-automorphic-seeds}

This chapter constructs the seed used in case II of the symplectic
support argument.  In particular, an automorphic representation is
produced, rather than added to the hypotheses of a deformation point.
The construction has three parts: an Eisenstein congruence with different
ordinary refinements at two places; a change of weight inside an ordinary
Hecke algebra; and an explicit integral basis for an induced
four-dimensional representation.  The last part is essential: abstract
automorphy of an induction would not specify the residual matrix entries
required by the support argument.

The source endpoint is \texttt{SeedIIAlt}.  Its premise already includes
an arbitrary original tower satisfying \texttt{Ok}, \texttt{Two}, and
\texttt{Deep}.  We preserve that tower, its labels, its chosen element
$g_0$, all its specified valuations, and the previously fixed field
isomorphism $\iota:\overline{\mathbb Q}_\ell\simeq\mathbb C$.
The only choices made afterwards are the seed and a finite enlargement
of the coefficient field.  There is no assertion here that an arbitrary
tower satisfies these three premises.

\subsection{Conventions and precisely stated classical inputs}

Write $\chi$ for the cyclotomic character, with arithmetic Frobenius
sent to the residue cardinality.  The source's attachment convention
uses geometric Frobenius: if $B_v$ is the normalized Satake multiset
of an automorphic representation of rank $n$, its attached Galois
multiset is
\[
                 q_v^{(n-1)/2} B_v.                 \tag{SX.1}
\]
Thus twisting a Galois representation by $\chi$ corresponds to
twisting its automorphic representation by $|\det|$.  This convention
accounts for the extra $|\det|$ in the rank-two to rank-four induction
below.

We use the following established theorems.  The applications and the
remaining algebra are proved in the chapter.

\begin{itemize}
\item \textbf{Hilbert Eisenstein series (SX-E).}
For a totally real field and finite ray characters $\eta,\zeta$ of
the appropriate parity, the holomorphic series $E_j(\eta,\zeta)$
has positive ideal coefficient
$\sum_{\mathfrak a\mid\mathfrak m}
 \eta(\mathfrak m/\mathfrak a)\zeta(\mathfrak a)
 N\mathfrak a^{j-1}$.
The weight-one series $E_1(1,\zeta)$ exists when $\zeta$ is totally
odd and nontrivial.  Its standard-cusp constants and its degeneracy
maps are those of this normalized expansion.
Dasgupta--Darmon--Pollack, Proposition 2.1 and its proof, printed
pp.~458--462, give this precise weight-one statement.  We also use
the usual decomposition of holomorphic Hilbert forms of weight $>2$
into cuspidal and Eisenstein eigensystems, and the rational
$q$-expansion principle.

\item \textbf{An integral Hasse invariant (SX-H).}
For a Hilbert modular variety at a prime unramified in the totally
real field, the determinant Hasse invariant has parallel weight
$\ell-1$, extends to the minimal compactification, and has expansion
$1$ at every Tate cusp.  A sufficiently high power lifts to
characteristic zero.  Mok--Tan, \S2.6, especially Proposition 2.11
on printed p.~906, gives the lifting by Serre vanishing; we use its
genus-one specialization.  Finite level descent by a product, described
in SX2, gives exactly the level-one form needed here.

\item \textbf{Ordinary control and ordinary Galois representations
(SX-O).}
At tame level prime to $\ell$, the cuspidal nearly ordinary Hecke
algebra with fixed algebraic central character is finite and
torsion-free over the corresponding weight algebra.  Its arithmetic
specializations of dominant algebraic weight are the classical
ordinary Hecke algebras after inverting $\ell$.  The identification
preserves away-from-$\ell$ Hecke operators and the normalized
operators $U_v^{\circ}$.  The attached rank-two Galois representation
has the usual unramified characteristic polynomials and the ordered
ordinary local characters specified by $U_v^{\circ}$ and the two
weight coordinates.  This statement applies to residual Eisenstein
systems.  Hida's \emph{Control of nearly ordinary Hecke algebras},
Theorems 1.2 and 2.4, Corollary 2.5, and Theorem 3.1, pp.~3--8,
state precisely these facts.  Remark 3.1 requires residual absolute
irreducibility only for matrices over the integral big Hecke algebra;
we only use individual characteristic-zero representations.

\item \textbf{Cyclic base change and automorphic induction (SX-BC).}
Cyclic base change for cuspidal $\mathrm{GL}_n$ representations and
cyclic automorphic induction exist, commute with local parameters,
and have the standard cuspidality criteria: a prime-degree base change
fails to be cuspidal precisely in the self-twist case, and an induction
of a cuspidal representation in degree two fails to be cuspidal
precisely when its two conjugates agree.  We use successive cyclic
base changes for an abelian extension.  Arthur--Clozel, Chapter III,
Theorems 4.2 and 5.1 and \S6, Theorem 6.2, are the global inputs.
For all-place compatibility of induction we use Henniart,
\emph{Induction automorphe globale pour les corps de nombres},
Theorems 4--5 and \S3.7, printed pp.~6--7 and 16.
Local compatibility here only concerns principal series and unramified
induction, where it is also the explicit character/norm description.

\item \textbf{Galois compatibility after CM base change (SX-G).}
A regular algebraic polarized cuspidal representation over a CM field
has an attached semisimple Galois representation, with semisimplified
local--global compatibility away from $\ell$.
We use BLGGT, \emph{Potential automorphy and change of weight},
Theorem 2.1.1, solely for this rank-two conclusion.
The quadratic character of the CM extension may be multiplied into
the polarization character to fix its real-place sign; its norm
pullback is trivial.  The comparison with the already constructed
ordinary representation is by Chebotarev and semisimplicity.
No automorphy lifting theorem is used to construct the seed.

\item \textbf{Classical local and global background (SX-C).}
We use global and local class field theory, Chebotarev density,
Dirichlet's unit theorem, the local newvector conductor theorem,
and the principal-series part of the local Langlands correspondence.
The latter identifies the conductor and determinant with those of
the Weil--Deligne parameter.  The elementary conductor and unitary
estimates needed from it are proved below.  The finite coefficient
capture theorem RV4 is already reconstructed and is used with an
actual continuous group homomorphism, with the matrix argument
spelled out in SX7.  For four characters, the absence of ratios
$|\cdot|^{\pm1}$ implies irreducibility of normalized induction
by Bernstein--Zelevinsky, Theorem 4.2, printed p.~454.
\end{itemize}

\subsection{The original quadratic data}

Fix an original datum $D$ over $O_c$ and a tower $T$ with
\[
 D.\texttt{Ctx2},\qquad T.\texttt{Ok},\qquad
 T.\texttt{Two},\qquad T.\texttt{Deep}.             \tag{SX.2}
\]
Let $F=T.F$, $K=T.K$, and $L=FK=T.L$ inside the original closure
$\Omega$.  We identify their absolute Galois groups with subgroups
of $G_{\mathbb Q}=\operatorname{Gal}(\Omega/\mathbb Q)$.
Put $t_0=T.g_0\in G_F$, and let $p_i$ be the four characters
$\texttt{chi}(D,T.\mathrm{pr},i)$ with the original labeling.
We use $\xi,\psi,\omega,\vartheta$ for their scalar extensions to
$\overline{\mathbb Q}_\ell$.  Thus
\[
 \alpha=\chi\omega^{-1}\xi,
 \quad \Psi=\chi^e\psi,
 \quad \vartheta=\alpha\Psi,
 \quad p_i p_{3-i}=\vartheta.                         \tag{SX.3}
\]
An overline denotes reduction modulo the maximal ideal, after any
needed finite coefficient extension.

\begin{proposition}[SX1: the mixed quadratic input on the original tower]
\label{claim:SX1}
There is a totally real quadratic field $F_0\subset\Omega$, already
specified by $T.\texttt{Two}$, split at $2$ and $\ell$.  Its sign
character $\varphi:G_{\mathbb Q}\to\{\pm1\}$ has the following
properties.  On $H=G_F\cap G_{F_0}$,
\[
 (\bar p_0,\bar p_1,\bar p_2,\bar p_3)
       =(\bar\xi,1,\bar\xi,1),                    \tag{SX.4}
\]
while $t_0\notin G_{F_0}$ and
\[
 \bar p_1(t_0)=-\bar p_3(t_0),\quad
 \bar p_2(t_0)=-\bar p_0(t_0),\quad
 \bar p_3(t_0)\in\{1,-1\},\quad
 \bar\xi(t_0^2)\ne1.                              \tag{SX.5}
\]
The character $\xi|_{G_{F_0}}$ is totally odd, has finite
prime-to-$\ell$ order, is unramified outside the primes above $2$,
and is nontrivial on inertia at a prime of norm $2$.  Its reduction
has nontrivial square somewhere and is nontrivial on each
decomposition group above $\ell$.
For every $k$ divisible by $\ell-1$,
\[
 \mu:=\chi^k\xi\varphi
       \quad\hbox{satisfies}\quad
 \bar\mu|_{G_F}=\bar\vartheta|_{G_F}.              \tag{SX.6}
\]
\end{proposition}
\begin{proof}
The field and splitting assertions are literal clauses of
$T.\texttt{Two.quad}$.  Its kernel in $G_F$ is the kernel of
$T.\mathrm{rb}(w_{\Psi})$; the latter has order at most two.
The character taking $1$ on $G_{F_0}$ and $-1$ on its other coset
is multiplicative because this subgroup has index two, and continuous
because it is open.

The two possible labelings are
\[
 (\alpha\Psi,\Psi,\alpha,1),\qquad
 (\alpha,1,\alpha\Psi,\Psi).                        \tag{SX.7}
\]
As $\bar\omega=\bar\chi$, one has $\bar\alpha=\bar\xi$.
On $H$, $\bar\Psi=1$, proving (SX.4).  If $t_0$ were in $H$,
the second and fourth residual eigenvalues at $t_0$ would coincide,
contrary to $T.\texttt{Ok.dist}$.  Thus $\bar\Psi(t_0)=-1$.
Substitution in (SX.7) proves the first three assertions of (SX.5).
If $\bar\xi(t_0)^2=1$, then $\bar\xi(t_0)=\pm1$, and again two
of the four displayed eigenvalues coincide.  This proves the last
assertion.

The finite order and the ramification of $\xi$ are the clauses
$\texttt{Two.ord}$ and $\texttt{Two.xi}$; its order is a power of
two, whereas $\ell$ is odd.  Its oddness follows from $D.\texttt{Good}$
and conjugacy of complex conjugations.  At either prime above $2$,
$\texttt{Two.big2}$ supplies inertia on which
$\bar\xi^4\ne1$; the splitting of $F_0$ at $2$ puts that inertia
inside $G_{F_0}$.  The square of this element proves that
$\bar\xi^2$ is not the trivial character.  Similarly
$\texttt{Two.bigL}$ and splitting at $\ell$ give the required
nontriviality on both decomposition groups above $\ell$.
In particular the conductor of $\xi|_{G_{F_0}}$ is nontrivial.

Finally $\bar\chi^k=1$, and on $G_F$ the characters $\bar\varphi$
and $\bar\Psi$ agree: both are $1$ on $H$ and $-1$ on its other
coset.  Hence $\bar\mu=\bar\xi\bar\Psi=\bar\vartheta$.
Every statement is on the original groups; no isomorphic copy of
the tower has been substituted.
\end{proof}

\subsection{An actual cuspidal ordinary congruence}

\begin{proposition}[SX2: mixed refinements force a cuspidal eigenlift]
\label{claim:SX2}
Let $F_0$ and $\xi$ be as in SX1, and fix $\iota$.
Choose one of the two primes above $\ell$, and call the singleton
set consisting of it $P$.  There exists an arbitrarily large
$k\ge2$, divisible by $\ell-1$, and a cuspidal ordinary Hilbert
eigenform of parallel weight $k+1$, nebentype corresponding to
$\xi$ in the geometric-Frobenius convention, and tame level its
conductor $\mathfrak c$, whose residual arithmetic Hecke system
is $1\oplus\bar\xi^{-1}$.  At the primes above $\ell$
the residual ordinary refinement is $\bar\xi^{-1}$ on $P$ and $1$
outside $P$.
\end{proposition}
\begin{proof}
Use class field theory and the fixed $\iota$ to write the finite
character as a complex narrow-ray character $\zeta$ with
$\zeta(v)=\iota(\xi(\operatorname{Frob}_{v,\mathrm{arith}})^{-1})$
at unramified places.  Its arithmetic Galois avatar is therefore
$\xi^{-1}$.  The series $E=E_1(1,\zeta)$
has first positive ideal coefficient $1$ and integral positive
coefficients.  At each $v\mid\ell$, its two stabilization roots
are $1$ and $\zeta(v)$.  Let $V_v$ be the degeneracy operator
whose normalized positive coefficient at $\mathfrak m$ is the
coefficient at $\mathfrak m/v$, or zero if $v\nmid\mathfrak m$.
Here is its constant-term normalization.  Let $u_v$ be the finite
idele equal to $\varpi_v^{-1}$ at $v$ and $1$ elsewhere.
For weight $j$, set
$V_v=q_v^{-j/2}R(\operatorname{diag}(u_v,1))$.
Substituting in the adelic Fourier expansion cancels the factor
$|u_v|^{j/2}=q_v^{j/2}$ and sends the normalized constant
function $c_0(y)$ to $c_0(yu_v)$.  The positive coefficient
indices shift by $v$, giving the stated formula.
For $E_1(1,\zeta)$, formula (50) of SX-E gives
$c_0(y)=2^{-[F_0:\mathbb Q]}L(\zeta,0)$, independently of
the narrow-class component.  Only that first constant term
occurs, since $\zeta$ has nontrivial conductor by SX1.
Apply
\[
 E_P=\prod_{v\in P}(1-V_v)
       \prod_{v\mid\ell,\,v\notin P}(1-\zeta(v)V_v)E.
                                                            \tag{SX.8}
\]
The identity $(U_v-1)(U_v-\zeta(v))E=0$, or the divisor-sum
coefficient formula, shows that $U_v E_P=\zeta(v)E_P$ on $P$
and $U_v E_P=E_P$ elsewhere.  At every standard cusp the constant
of $V_vE$ equals that of $E$ in this normalization.  Since $P$
is nonempty, the first product in (SX.8) kills these constants.
All positive coefficients remain integral, and the first one remains
$1$.  This assertion concerns the standard expansions on every
narrow-class component.  It does not assert that every cusp constant
of $E_P$ vanishes, and does not assert that $E_P$ is cuspidal.

Here is the required form congruent to $1$.  Choose an auxiliary
prime-to-$\ell$ fine level.  The Hilbert determinant Hasse invariant
on its minimal compactification has parallel weight $\ell-1$.
On a Tate object the Cartier operator fixes the logarithmic
differentials, so its determinant has expansion $1$.  For all large
even $a$, Serre vanishing applied to multiplication by $\ell$ in the
Hodge line bundle shows that its $a$th power lifts to an integral
characteristic-zero section $G_a$ of weight $a(\ell-1)$.
Make the auxiliary cover Galois and multiply all deck translates
of $G_a$.  This product descends to level one, has trivial
nebentype, and has expansion $1$ modulo the maximal ideal at
every cusp.  This descent uses a product, so no division by the
order of the deck group is involved.  On all narrow-class components
the same canonical Hasse section is used.  To include the adelic
central descent, take the product also over the finite central
ideal-class action on these components.  The normalized line-bundle
action factors through this finite class group: principal positive
ideles act by the prescribed parallel weight, and integral units
by the level-one trivial character.  Choose ideal-class representatives
prime to $\ell$.  Because the starting weight is a multiple of
$2(\ell-1)$, the rational norm factors in the normalized central
action are $\ell$-adic units congruent to $1$.
The Hasse section is fixed
by this action, by the Cartier functoriality under prime-to-$\ell$
isogenies (Mok--Tan, Lemma 2.10).  Every factor therefore still
has reduction $1$.  Taking a further common power if needed
trivializes the finite stabilizer action on the Hodge line.
Invariant sections descend on the finite quotient, and the
resulting adelic form has the trivial finite central character.
Increasing $a$ gives
arbitrarily large resulting weights $k$ divisible by $\ell-1$.

Multiply $E_P$ by such a level-one form $G$.  In a standard expansion
the positive coefficient of the product is the constant coefficient
of $G$ times that of $E_P$, plus a sum of products of positive
coefficients.  The possibly nonintegral constant of $E_P$ is zero.
Consequently every positive coefficient of $E_PG$ is integral,
and its reduction equals that of $E_P$.  The first coefficient is a
unit.  The explicit Hecke coefficient formulas then show that
$E_PG$ is an approximate simultaneous eigenvector modulo the
maximal ideal for all the away Hecke operators and all $U_v$.
At weight $k+1$ the factor $(Nv)^k$ occurring in these formulas
reduces to $1$ for $v\nmid\ell$, as $\ell-1\mid k$.

For completeness, the eigenlift is an algebraic step.  Work over a
finite coefficient extension containing $E_PG$, its character, and
the rational holomorphic space.  The module of forms with integral
positive coefficients is a lattice: finitely many coefficient
functionals already inject the finite-dimensional space into a
finite product of the coefficient field, and the $q$-expansion
principle supplies bounded denominators in the other direction.
It is Hecke stable by the coefficient formulas.  Let $\mathbb T$
be the finite commuting Hecke algebra acting on it.  The nonzero
residual approximate eigenvector gives a maximal ideal
$\mathfrak m\subset\mathbb T$ with the specified residue
eigenvalues.  The localized lattice is nonzero and torsion-free
over the discrete valuation ring.  Hence some prime of
$\mathbb T_{\mathfrak m}$ avoids $\ell$.  After passing to a
finite characteristic-zero residue field and a common eigenline,
one obtains a holomorphic eigenform with exactly this residual
system.  Its $U_v$ eigenvalues are units, so it is ordinary.

It remains to prove that this lift is cuspidal.  Suppose otherwise.
Its Eisenstein system has two finite characters $\eta_1,\eta_2$,
with away eigenvalues
$\eta_1(v)+(Nv)^k\eta_2(v)$ and product
$\eta_1\eta_2=\zeta$.  Reduction and Chebotarev give, for every
Galois element $g$,
\[
 \bar\eta_1(g)+\bar\eta_2(g)=1+\bar\xi(g)^{-1},
 \qquad\bar\eta_1(g)\bar\eta_2(g)=\bar\xi(g)^{-1}.
                                                            \tag{SX.9}
\]
Thus $(\bar\eta_1(g)-1)(\bar\eta_2(g)-1)=0$.
A group cannot be the union of two proper subgroups: if $a$ is
outside the first kernel and $b$ outside the second, each lies in
the other, whereas $ab$ lies in neither.  Hence one of the two
characters is identically $1$ and the other is $\bar\xi^{-1}$.

At $v\mid\ell$, both inducing characters are unramified.  Indeed
the nebentype is unramified there and the level has exponent one;
if either character were ramified, cancellation in the central
character would make both ramified and their conductor sum would
be at least two.  The two possible $U_v$ eigenvalues are
$\eta_1(v)$ and $(Nv)^k\eta_2(v)$.  Only the first is a unit.
Therefore the residual ordinary eigencharacter is the same
global character $\bar\eta_1$ at every $v\mid\ell$.
Our prescribed refinement is $\bar\xi^{-1}$ at the prime in $P$
and $1$ at the other prime, while $\bar\xi^{-1}$ is nontrivial at both.
Neither of the possibilities following (SX.9) is compatible with
these two refinements.  This contradiction proves cuspidality.
It is this argument, rather than multiplication by $G$, that kills
the Eisenstein alternative.
\end{proof}

\subsection{Changing weights with fixed determinant}

\begin{proposition}[SX3: the full mixed rank-two representation]
\label{claim:SX3}
On the same $F_0$, $\xi$, and fixed $\iota$, there are integers
$k,t\ge2$, each divisible by $\ell-1$, and a continuous
representation $r:G_{F_0}\to\mathrm{GL}_2(\overline{\mathbb Q}_\ell)$
with the following properties:
\[
 \det r=\chi^k\xi,
 \qquad\bar r^{\mathrm{ss}}=1\oplus\bar\xi.         \tag{SX.10}
\]
It is irreducible on every open subgroup.  Away from $\ell$,
its inertia is conjugate to $\operatorname{diag}(1,\xi)$.
At the two places above $\ell$ it has ordered ordinary flags
with inertia weights and residual diagonal characters
\[
 \begin{array}{c|c|c}
    &\text{inertia weights}&\text{residual ordered characters}\\ \hline
   v\in P &(k+t,-t)&(\bar\xi,1)\\
   v\notin P &(k,0)&(1,\bar\xi).
 \end{array}                                                   \tag{SX.11}
\]
There is an actual regular algebraic cuspidal representation
$\pi_0$ of $\mathrm{GL}_2(\mathbb A_{F_0})$ attached to $r$ in
(SX.1), with infinitesimal weights
$\frac12+(k+t,-t)$ and $\frac12+(k,0)$.
Every finite component of $\pi_0$ is a principal series with no
linked pair of inducing characters.  Thus the literal source
$\texttt{MixDat}$ and $\texttt{PR}$ conclusions are supplied.
\end{proposition}
\begin{proof}
Start with SX2 and its residual ordinary maximal ideal.
Fix the algebraic central character, including its norm power;
equivalently, fix the determinant $\chi^k\xi^{-1}$ of the usual
covariant Hilbert Galois representation.  In the notation of
Hida's control theorem the weight coordinates are initially
$(\kappa_1,\kappa_2)=(0,k-1)$ at both embeddings.  They may vary
with their sum fixed at $k-1$.  This parallel sum is exactly the
condition imposed by totally positive central units.
Here $F_0$ is real quadratic, so one may use the quaternion
algebra ramified at its two real places and split at every finite
place.  Jacquet--Langlands preserves the finite Hecke operators
and the prescribed tame level.  Its definite ordinary module
satisfies the finite projectivity and control statements of
SX-O (Hida, Lemmas 2.1--2.2); passing to sufficiently deep
$\ell$-level removes stabilizers.  The descent in Theorem 1.2
uses faithful flatness of weight algebras, so it does not require
dividing by a tame-level stabilizer order.
Localize the cuspidal
ordinary Hecke algebra at the maximal ideal of SX2 and take its
completed nonzero factor $B$.  The relevant component of the
weight algebra $\Lambda$ is a power-series domain.  By SX-O,
$B$ is finite and torsion-free over $\Lambda$.
In particular $\Lambda\to B$ is injective.

Choose $t\ge2$ divisible by $\ell-1$ and, if necessary, by an
arbitrarily high power of $\ell$.  Specialize the weights to
$(\kappa_1,\kappa_2)=(-t,k+t-1)$ at $P$ and $(0,k-1)$ outside $P$.
They have the required central sum and the same torsion character
modulo the maximal ideal.  Integral lying-over for
$\Lambda\hookrightarrow B$ gives a prime over this arithmetic
weight which avoids $\ell$; a finite extension of the specialization
field contains its residue field.  The control isogeny identifies
this characteristic-zero point with a classical cuspidal nearly
ordinary eigenform.  It has the original tame level and central
character.  The normalized $U_v^\circ$ and the away Hecke
operators retain the residual values from SX2.  This proves the
existence of the new form at the required nonparallel weights;
it is not an assumption about a hypothetical component through a
Galois deformation point.

Let $r_f$ be the usual covariant Galois representation of this
form.  Hida's Theorem 3.1 writes its ordinary flag as
$\left(\begin{smallmatrix}\beta&*\\0&\alpha\end{smallmatrix}\right)$,
where $\alpha$ has inertia $\chi^{\kappa_1}$ and $\beta$ has
inertia $\chi^{\kappa_2+1}$.  The normalized unit eigenvalue is
that of the \emph{quotient} $\alpha$ after removing its algebraic
inertia weight.  It has residual character $\bar\xi^{-1}$ on
$P$ and $1$ elsewhere.  These assertions hold on the entire
decomposition group: they agree on inertia and on an unramified
Frobenius lift, which together topologically generate it.
Define
\[
                         r=\chi^k r_f^\vee.
\]
The annihilator of the ordinary subline is the subline in the
dual.  Its character is $\alpha^{-1}$, so the upper character
of $r$ has weight $k+t$ and residual character $\bar\xi$ on
$P$, and weight $k$ and residual character $1$ elsewhere.
The lower characters are determined by the determinant.  Thus
the ordered flags are exactly (SX.11), with no reversal hidden
in the ordinary terminology.  Good-prime characteristic
polynomials and Chebotarev give
$\det r_f=\chi^k\xi^{-1}$ and
$\bar r_f^{\mathrm{ss}}=1\oplus\bar\xi^{-1}$.
Dualizing and twisting, with $\bar\chi^k=1$, proves (SX.10).
All characters are continuous; their finite-image parts have open
kernels.  The ordinary theorem is being used only at the
characteristic-zero specialization, so no assumption of residual
irreducibility is present.

Here are two elementary irreducibility checks.  If $r$ had a
global invariant line, its character would be unramified outside
finitely many places and would have one of the weights $k+t,-t$
at the first embedding and one of $k,0$ at the second.  A global
character of a totally real field with such algebraic local
inertia has parallel weights.  To see this directly, apply class
field theory to totally positive global units sufficiently close
to $1$ at the finitely many ramified places.  The product of their
embedding powers, with these integer weights, is $1$.  Apply the
fixed $\iota$ and the real logarithm.  Dirichlet's unit lattice
spans the hyperplane of vectors with coordinate sum zero, so the
weights must all be equal.  The two sets here are disjoint because
$k,t\ge2$.  This proves absolute irreducibility.

Suppose next that restriction to an open subgroup were reducible,
and replace that subgroup by its open normal core $N$.
Inertia above $\ell$ has open cyclotomic image, even after
intersection with $N$; hence $r(N)$ contains an element with two
distinct eigenvalues, using either nonzero weight difference in
(SX.11).  A line invariant under $N$ is therefore one of its two
eigenlines.  Normality makes the translates of that line again
$N$-invariant, so the full group permutes these two lines.
Irreducibility makes the permutation nontrivial.
Let $J$ be the index-two subgroup fixing each line.  Elements
outside $J$ have trace zero, hence reduction of (SX.10) gives
$\bar\xi=-1$ on that coset.  Multiplication by any element of
$J$ then gives $\bar\xi=1$ on $J$.  Thus $\bar\xi^2=1$
everywhere, contradicting SX1.  The restriction to every open
subgroup is irreducible.

We give the away local-type comparison carefully.  At a place
outside $\ell\mathfrak c$, the ordinary Galois theorem already
gives unramifiedness.  At a prime dividing $\mathfrak c$, the
automorphic conductor is at most $a(\xi)$, and its central
character has conductor exactly $a(\xi)$.  Thus both conductors
are equal.  A two-dimensional Weil--Deligne parameter with this
equality is a direct sum of two characters, exactly one
unramified, with zero monodromy.  Indeed an irreducible Weil
parameter has no inertia invariants: its tame conductor term is
two, whereas that of its determinant is at most one, and on
every higher ramification group determinant codimension is at
most vector codimension.  The conductor inequality is then
strict.  A special parameter also has strictly larger conductor
than its determinant.  Finally for two ramified characters the
conductor sum is strictly larger than the conductor of their
product.  These observations prove the assertion.  The remaining
ramified character restricts to $\xi$ or its inverse according
to the contravariant normalization.

To transfer this assertion to the actual Galois representation,
choose an imaginary quadratic extension of $\mathbb Q$ splitting
at the rational prime under the place being inspected.  Its
compositum $E/F_0$ is CM quadratic and splits at that place.
Such a quadratic field exists by the elementary congruence
conditions on a negative squarefree discriminant.
The cyclic base change of the Hilbert cusp is cuspidal: a
nontrivial self-twist would give a reducible restriction of $r$
to this open subgroup, just excluded.  Its central character
provides the polarization.  Multiplying the polarization character
by the odd quadratic character of $E/F_0$, if required, changes
its real-place sign and leaves its norm pullback unchanged.
SX-G applies.  Its representation equals $r|_{G_E}$ by good
Frobenius characteristic polynomials and semisimplicity.
The split local group is the original one.  Semisimplified local
compatibility therefore gives the inertial characters $1,\xi$.
They are distinct at a ramified conductor place; monodromy
commutes with inertia and hence is diagonal, so its nilpotence
forces it to vanish.  Thus the actual inertia representation,
not just a list of Frobenius traces, is
$\operatorname{diag}(1,\xi)$.  This argument avoids imposing
an unavailable finite discrete-series place in an even-degree
Hilbert local--global theorem.

At $v\mid\ell$ the classical form has Iwahori level and
unramified central character.  Its local component is either
unramified principal series or an unramified twist of Steinberg.
The Steinberg possibility has normalized ordinary slope
$(k-1)/2$ outside $P$ and $(k-1)/2+t$ on $P$, both positive.
It is excluded by the unit $U_v^\circ$ eigenvalue.  Thus it is
unramified principal series there as well.

To specify the attachment normalization, let $\pi_f$ be the
unitary cuspidal representation of this form and put
\[
              \pi_0=\pi_f\otimes
                         |\det|^{(k+1)/2}.          \tag{SX.12}
\]
If $\lambda_1,\lambda_2$ are the arithmetic eigenvalues of
$r_f$ at an unramified place, the Satake parameters of
$\pi_f$ are $q^{-k/2}\lambda_i$.  The geometric eigenvalues
of $r=\chi^k r_f^\vee$ are $q^{-k}\lambda_i$.
The parameters of (SX.12) are
$q^{-(k+1)/2}q^{-k/2}\lambda_i=q^{-k-1/2}\lambda_i$.
Multiplication by $q^{1/2}$ therefore gives exactly (SX.1).
The infinitesimal characters are $\frac12+(k+t,-t)$ and
$\frac12+(k,0)$, proving regular algebraicity.

We also record a local estimate needed in SX6.  For a unitary
irreducible principal series of $\mathrm{GL}_2$ with one
unramified and one ramified inducing character, both characters
are unitary: comparison with its conjugate dual either fixes
each character or exchanges them, and the exchange would force
their restrictions to units to agree.  At a spherical place write
the Satake parameters as $a,b$.  Conjugate duality and the unitary
central character imply either $|a|=|b|=1$, or
$(a,b)=(zu,z^{-1}u)$ with $z>1$ and $|u|=1$.
On the unit spherical vector, the normalized Hecke operator is
$q^{-1/2}$ times a sum of $q+1$ unitary translates.  Therefore
\[
                   |a+b|\le q^{1/2}+q^{-1/2}.
\]
It follows that $z\le q^{1/2}$.  Equality would force all these
translates to span the same line, by equality in the Hilbert-space
triangle inequality.  That line, fixed by the maximal compact
subgroup and stable under its uniformizer double coset, is stable
under $\mathrm{GL}_2$.  This would make the representation
one-dimensional, impossible for a cuspidal local component.
Thus $q^{-1/2}<|a|,|b|<q^{1/2}$.
After (SX.12), all finite inducing characters $\nu$ satisfy
\[
 q_v^{s-1/2}<|\nu(\varpi_v)|<q_v^{s+1/2},
                  \qquad s=-(k+1)/2.               \tag{SX.13}
\]
The same inequalities exclude a quotient $\nu_1/\nu_2=|\cdot|^{\pm1}$,
so these are the actual unlinked principal series required by
$\texttt{PlainR}$.
\end{proof}

\subsection{An integral induced representation with specified entries}

\begin{proposition}[SX4: the normalized Ribet lattice and mixed induction]
\label{claim:SX4}
After conjugating $r$, there is a continuous representation
$R:G_{\mathbb Q}\to\mathrm{GSp}_4(\overline{\mathbb Q}_\ell)$
with multiplier $\mu=\chi^k\xi\varphi$, abstractly
$\operatorname{Ind}_{G_{F_0}}^{G_{\mathbb Q}}r$, such that its
restriction to $G_F$ is integral, its diagonal reduces to
$(\bar p_i)$, and its off-diagonal residual support is contained
in
\[
        vV=\{w_\beta,-w_\alpha,-w_\vartheta\}.
                                                            \tag{SX.14}
\]
It has a unit entry of weight $-w_\alpha$, or it has both a unit
entry of weight $w_\beta$ and a unit entry of weight
$-w_\vartheta$.  Also $R(t_0)$ is diagonal.
\end{proposition}
\begin{proof}
Let $u=t_0^2\in H$.  Its two residual eigenvalues are distinct
by (SX.5).  Diagonalize $r(u)$ with its first eigenvalue reducing
to $1$ and its second to $\xi(u)$.  For $h\in H$, compare
$\operatorname{tr}r(h)$ and $\operatorname{tr}r(uh)$ with their
residual values.  The determinant of the resulting two-by-two
system is the unit difference of the two eigenvalues of $r(u)$.
It follows that
\[
       r(h)_{00}\equiv1,\qquad r(h)_{11}\equiv\xi(h).
                                                            \tag{SX.15}
\]
Comparing the $(0,0)$ entry of $r(hh')=r(h)r(h')$ then gives
$|r(h)_{01}r(h')_{10}|<1$ for every $h,h'\in H$.
Since $r|_H$ is irreducible, its upper-right entry is not
identically zero.  Compactness of $H$ and continuity give a
largest absolute value of that entry, attained at some $h_*$.
Conjugation by a diagonal matrix makes this maximum $1$.
The preceding product inequality now gives
\[
 r(h)_{00}\equiv1,\quad r(h)_{11}\equiv\xi(h),\quad
 |r(h)_{10}|<1,\quad |r(h)_{01}|\le1,
 \qquad |r(h_*)_{01}|=1.                            \tag{SX.16}
\]
The diagonalization of $r(u)$ is preserved.  This is an explicit
lattice normalization; it requires no assertion that an arbitrary
compact subset of $\overline{\mathbb Q}_\ell$ lies in a finite field.

Choose a square root $a$ of $r(u)_{00}$ reducing to $p_3(t_0)$,
and put $b=\mu(t_0)/a$.  The two choices of square root are
distinguished because $\ell\ne2$.  Equations (SX.5)--(SX.6)
give $\bar b=\bar p_0(t_0)$, and the determinant identity gives
$r(u)=\operatorname{diag}(a^2,b^2)$.
Both $a$ and $b$ are units.  Set $\gamma=b/a$.

In the usual induced basis, $h\in G_{F_0}$ acts by
$\operatorname{diag}(A,B)$, where
$A=r(h)$ and $B=r(t_0^{-1}ht_0)$, and $t_0$ acts by
$\left(\begin{smallmatrix}0&\operatorname{diag}(a^2,b^2)\\1&0\end{smallmatrix}\right)$.
Change basis by
\[
 P=\begin{pmatrix}
 0&a&0&a\\ b&0&-b&0\\0&-1&0&1\\1&0&1&0
 \end{pmatrix}.
                                                            \tag{SX.17}
\]
Direct multiplication gives $R(t_0)=\operatorname{diag}(b,-a,-b,a)$
and
\[
R(h)=\frac12\begin{pmatrix}
A_{11}+B_{11}&A_{10}/\gamma-B_{10}&B_{11}-A_{11}&A_{10}/\gamma+B_{10}\\
\gamma A_{01}-B_{01}&A_{00}+B_{00}&-\gamma A_{01}-B_{01}&A_{00}-B_{00}\\
B_{11}-A_{11}&-A_{10}/\gamma-B_{10}&A_{11}+B_{11}&-A_{10}/\gamma+B_{10}\\
\gamma A_{01}+B_{01}&A_{00}-B_{00}&-\gamma A_{01}+B_{01}&A_{00}+B_{00}
\end{pmatrix}.                                                \tag{SX.18}
\]
The group law follows from the induced representation and this
one fixed basis change.  For the multiplier, use the alternating
form $\operatorname{diag}(J_2,abJ_2)$ on the induced space, where
$J_2=\left(\begin{smallmatrix}0&1\\-1&0\end{smallmatrix}\right)$.
Its Gram matrix in (SX.17) is $-2abJ$, with $J$ the original
four-dimensional alternating matrix.  Both the subgroup and
$t_0$ act on it with multiplier $\mu$.  This proves the exact
symplectic identity, including the quadratic sign $\varphi$.

For $h\in H$, both $A$ and $B$ satisfy (SX.16), since $H$ is
normal in $G_F$ and $\xi$ is a global character.  Formula (SX.18)
therefore is integral, reduces on the diagonal to
$(\bar\xi,1,\bar\xi,1)$, and has possible nonzero off-diagonal
entries only at $(1,0),(1,2),(3,0),(3,2)$.
Their weights are respectively
$-w_\alpha,w_\beta,-w_\vartheta,-w_\alpha$.
All other entries reduce to zero by (SX.16).
At $h_*$, put $x=\gamma A_{01}$ and $y=B_{01}$.
The $(1,0)$ and $(3,0)$ entries are $(x-y)/2$ and $(x+y)/2$;
their sum is the unit $x$.  Hence one is a unit.
In the second case the $(1,2)$ entry is the negative of the
$(3,0)$ entry, so it too is a unit.
This proves precisely the asserted alternative.

Every remaining element of $G_F$ is an element of $H$ times
$t_0$.  Right multiplication by the diagonal unit matrix
$R(t_0)$ preserves the support and integrality.  Its residual
diagonal is correct by (SX.5) and the choices of $a,b$.
Continuity follows from the two open cosets and the explicit
formulas.  This proves all assertions on the original group.
\end{proof}

\subsection{The actual local flags and the correcting character}

\begin{proposition}[SX5: local flags, inertia kernel, and the exact weights]
\label{claim:SX5}
There is a continuous character $\delta:G_F\to1+\mathfrak m_{\bar{\mathbb Q}_\ell}$
and, for each original index $v\in T.\iota$, an integral lower
unipotent symplectic matrix $n_v$ such that
\[
        Q_{\bar{\mathbb Q}_\ell}(g)=\delta(g)R(g),
        \qquad\delta^2\mu=\vartheta.              \tag{SX.19}
\]
These matrices satisfy every global and local equation of
$(T.\mathrm{fr}(\varnothing)).\texttt{IsPair}$.
Their residual local support is contained in
$qQ=\{-w_\alpha,-w_\Psi,-w_\vartheta\}$ and
$(n_v)_{20}$ is a unit.  Put
\[
 m=(k+t,k,0,-t),\qquad b_0=e+1-k,\qquad h_i=2m_i+b_0.
                                                            \tag{SX.20}
\]
Then the actual pair has source weights $h_i$, with
$m_i-m_{i+1}\ge2$.  There is one integer $N>0$ such that
\[
                (\delta(g)^2\chi(g)^{-b_0})^N=1
                         \quad(g\in G_F).          \tag{SX.21}
\]
\end{proposition}
\begin{proof}
The character $\vartheta/\mu$ is continuous and reduces to $1$
by SX1.  Squaring is an automorphism of the group of units
congruent to $1$ when $\ell$ is odd.  Its inverse is continuous:
on any finite coefficient field it is the unique Hensel root
congruent to $1$, or equivalently the convergent binomial series.
Uniqueness makes the pointwise root multiplicative.  This defines
$\delta$, proves (SX.19), and preserves every residual entry
from SX4.

Fix an original valuation above $\ell$.  It splits in $F_0$,
so its decomposition group lies in $G_{F_0}$; conjugation by
$t_0$ exchanges the two places.  At the one whose top residual
character is $\xi$, a common ordinary eigenline can be written
as $(u,1)$ with $|u|\le1$.  At the other it can be written as
$(1,w)$ with $|w|<1$.  To check these assertions, use an element
of the local group for which $\bar\xi\ne1$, supplied by SX1.
In the eigenvector equation its eigenvalue minus the opposite
diagonal entry is a unit.  Formula (SX.16) then bounds the ratio
of the two coordinates by the upper-right integral entry in the
first case and by the strictly small lower-left entry in the
second.  The required denominator coordinate cannot be zero.
The same line works for the whole local group because it is the
ordinary invariant line.

For $n(p,q,r,s,z)$ write
\[
 n(p,q,r,s,z)=
 \begin{pmatrix}1&0&0&0\\p&1&0&0\\q&s&1&0\\r&z&-p&1\end{pmatrix}.
                                                            \tag{SX.22}
\]
It is symplectic exactly when $z+ps-q=0$.
If the first block in (SX.18) has eigenline $(u,1)$, choose
\[
 n_v=n\left(\gamma u,-1,\gamma u,
       \frac{-2w}{1+\gamma uw},
       \frac{-(1-\gamma uw)}{1+\gamma uw}\right).     \tag{SX.23}
\]
If that eigenline occurs in the second block, choose instead
\[
 n_v=n\left(-u,1,u,
       \frac{-2w}{\gamma+uw},
       \frac{\gamma-uw}{\gamma+uw}\right).           \tag{SX.24}
\]
The denominators are units because $|\gamma|=1$ and $|uw|<1$.
These formulas satisfy $z+ps-q=0$, all their entries are
integral, their $(2,1)$ entry is strictly small, and their
$(2,0)$ entry is $-1$ or $1$.
This proves precisely the residual support and nonzero entry
assertions.

For clarity, triangularization in (SX.23) is an identity in six
scalars, not a dimension argument.  If $A(u,1)^t=a_1(u,1)^t$,
$\det A=a_1a_2$, then
\[
 A=\begin{pmatrix}
 a_2+ux&(a_1-a_2-ux)u\\x&a_1-xu
 \end{pmatrix}.
\]
If $B(1,w)^t=b_1(1,w)^t$, $\det B=b_1b_2$, then
\[
 B=\begin{pmatrix}
 b_1-yw&y\\(b_1-b_2-yw)w&b_2+yw
 \end{pmatrix}.
\]
Substitution into (SX.18) and (SX.23) makes
$n_v^{-1}R(g)n_v$ upper triangular.  Its diagonal is
$(a_1,b_1,b_2,a_2)$.
Interchanging the blocks gives (SX.24), whose diagonal is
$(b_1,a_1,a_2,b_2)$.
In both cases (SX.11) gives the inertia exponents
$m=(k+t,k,0,-t)$ and residual diagonal
$(\bar\xi,1,\bar\xi,1)$.
Thus the original local character order is respected.

We next verify the global inertia kernel in the frame.
Let $g\in G_F$ be an inertia element away from $\ell$ with
$D.\bar r(g)=1$.  The clause $\texttt{Two.unr}$ puts it in
$G_{F_0}$.  Away from $\ell$ one has $\chi(g)=\omega(g)=1$;
the residual diagonal then gives $\bar\xi(g)=1$.
The order of $\xi$ is prime to $\ell$, so $\xi(g)=1$.
Also $\texttt{Ok.kb}$ gives $\psi(g)=1$.
By SX3 both $r(g)$ and its $t_0$-conjugate are the identity.
Hence $R(g)=1$ by (SX.18), and (SX.19) gives $\delta(g)^2=1$.
The root congruent to $1$ is $1$, so $Q_{\bar{\mathbb Q}_\ell}(g)=1$.
The multiplier, diagonal-at-$g_0$, and local flag clauses were
already proved in (SX.19), SX4, and (SX.23)--(SX.24).

Dividing the two multipliers in (SX.19) gives the exact identity
\[
       \delta^2=\chi^{e+1-k}\omega^{-1}\psi\varphi^{-1}
                                  \quad\hbox{on }G_F.       \tag{SX.25}
\]
Choose $a>0$ prime to $\ell$ annihilating $\psi|_{G_F}$, as
provided by $\texttt{Deep}$.  We may take
$N=2a(\ell-1)$ in (SX.21).
On the open subgroup $U=\ker\omega\cap\ker\psi$ of $G_F$,
the inertia at every place above $\ell$ also kills $\varphi$
by splitting, so $\delta^2=\chi^{b_0}$ there.
The local diagonal characters $\delta\chi^{m_i}$ therefore
have square $\chi^{2m_i+b_0}$, exactly the source's
$\texttt{Wts}$ convention.

Finally the source's literal set in the local frame is
\[
 \operatorname{muOf}(D_v)=
 \{s\in D_v:\ s^{\ell^a}\in
       \overline{[D_v,D_v]}\text{ for some }a\ge0\}.
\]
It is contained in $I_v\cap\ker\chi$; equality is not needed.
Indeed $D_v/I_v\simeq\widehat{\mathbb Z}$ is abelian and has
no torsion, so the indicated power relation puts $s$ in $I_v$.
The continuous abelian character $\chi$ kills the closed
commutator subgroup, giving $\chi(s)^{\ell^a}=1$.
For odd $\ell$, $\mathbb Z_\ell^\times$ has no nontrivial
$\ell$-power torsion: its torsion has order dividing $\ell-1$
and its principal units are torsion-free, by the $\ell$-adic
logarithm.  Hence $\chi(s)=1$.
For every such $s$, one has $\chi(s)=\omega(s)=1$.  The clause
$\texttt{Ok.ineL}$ makes its residual diagonal trivial, hence
$\bar\psi=1$.  The prime-to-$\ell$ order supplied by
$\texttt{Deep}$ implies $\psi=1$ on this subgroup.
Equation (SX.25) gives $\delta=1$, and all diagonal local
characters are $1$.  This is the exact frame condition, not just
an equality on a smaller chosen inertia subgroup.
The adjacent differences of $m$ are $t,k,t$, all at least two.
\end{proof}

\subsection{Irreducibility and the actual automorphic witness}

\begin{proposition}[SX6: cyclotomic irreducibility and the local window]
\label{claim:SX6}
The representation $R$ is irreducible on
$G_L\cap\ker\chi$.  There is a cuspidal regular algebraic
representation $\Pi$ of $\mathrm{GL}_4(\mathbb A_L)$ attached
to $R|_{G_L}$ in convention (SX.1).  It satisfies the literal
\texttt{ImaginaryQuadraticEtaCharacter.\allowbreak AllWin} condition:
there is a single real number $s_4$ such that every finite local
component is toral and every character $\nu$ in its cuspidal
support satisfies
\[
 q_v^{s_4-1/2}<|\nu(\varpi_v)|<q_v^{s_4+1/2}.
                                                            \tag{SX.26}
\]
One can take $s_4=-(k+3)/2$.
\end{proposition}
\begin{proof}
Put $V=G_L$ and $U=G_L\cap\ker\chi$.
Write $\eta=\bar\Psi|_{G_F}$ for the quadratic character
cutting out $F_0$.  We prove directly that $\eta|_U\ne1$.
Suppose $\eta|_U=1$.  If $g\in V$ and $\omega(g)=1$,
then $\chi(g)\equiv1\pmod\ell$, and
$\chi(g^{\ell^n})\longrightarrow1$.
Compactness gives a cluster point $h\in V$ of the sequence
$g^{\ell^n}$.  The character $\eta$ is quadratic and $\ell$
is odd, so each term lies in the same closed $\eta$-coset
as $g$.  Thus $\chi(h)=1$ and $\eta(h)=\eta(g)$.
The assumption on $U$ now gives $\eta(g)=1$.
Consequently $\eta|_V$ is quadratic and trivial on
$V\cap\ker\omega$.  Multiplication by $\eta$ permutes the
four residual characters: in (SX.7) it exchanges indices
$0\leftrightarrow2$ and $1\leftrightarrow3$.
The precise no-self-twist clause $\texttt{Ok.g3'}$ therefore
forces $\eta|_V=1$.  But $\texttt{Ok.g3}$ applied to
$w_\Psi$ asserts that it is nontrivial on $V$.
This contradiction proves the required nontriviality.
Thus $U_0=U\cap G_{F_0}$ has index two in $U$.

The restriction $r|_{U_0}$ is irreducible.  Here is the precise
infinite-subgroup argument.  The quotient of the open group
$G_{F_0}\cap G_L$ by $U_0$ is a subgroup of
$\mathbb Z_\ell^\times$, hence procyclic since $\ell$ is odd.
The finite-dimensional algebra spanned by the restricted image
has radical stable under the whole open group.  Its radical
annihilates a nonzero subspace and hence, by irreducibility of
that open-group representation, annihilates the entire space.
The restriction is semisimple.  Its finitely many isotypic
summands have an open stabilizer; more than one summand would
contradict SX3.  A repeated single constituent has a projective
multiplicity representation of the procyclic quotient.
A lift of a topological generator has an eigenline; polynomial
identities and continuity make that line invariant under the
whole quotient.  Irreducibility forces multiplicity one.
This is the same finite-algebra/procyclic argument proved in RX4.

The two $U_0$ constituents $r$ and $r^{t_0}$ are not isomorphic.
Suppose they were.  Their determinants agree on the surrounding
open group, since $\det r$ is the restriction of a global
character.  Fix an intertwiner on $U_0$.  Schur's lemma shows
that, for every element $x$ of this surrounding group, the
comparison of the two extensions is scalar.  Its square is
$1$ by the determinant equality.  Consequently the two
representations have the same characteristic polynomial at
$x^2$.  At a place above $\ell$, take an inertia element in
this open group with cyclotomic value of infinite order.
Such elements exist because the local cyclotomic image has
finite index in $\mathbb Z_\ell^\times$.
The two polynomials, evaluated at a suitable positive even
power, have eigenvalues
\[
     \{\chi(x)^{k+t},\chi(x)^{-t}\},\qquad
     \{\chi(x)^k,1\}.
\]
The second contains $1$ and the first does not, since
$t,k+t>0$.  This is a contradiction.  The usual index-two
induction argument now proves $R|_U$ irreducible: its restriction
to $U_0$ has two nonisomorphic irreducible constituents exchanged
by the other coset, so no proper invariant sum is possible.

In particular $F_0$ is not contained in $L$.  Put
$E=F_0K$ and $M=F_0L$.  Then $E/F_0$ is CM quadratic,
$M/E$ is abelian, and $M/L$ is quadratic.
Apply cyclic base change successively to $\pi_0$ to obtain
$\pi_M$.  All intermediate base changes are cuspidal:
the self-twist obstruction would contradict irreducibility of
$r$ on the corresponding open subgroup.  The representations
$\pi_M$ and its conjugate over $L$ are distinct, since otherwise
their good Frobenius polynomials would identify $r$ and its
conjugate on $G_M$, contrary to the preceding argument on
$U_0$.  Automorphic induction therefore gives a cuspidal
$\operatorname{AI}_{M/L}(\pi_M)$.
Set
\[
             \Pi=\operatorname{AI}_{M/L}(\pi_M)
                                      \otimes|\det|.          \tag{SX.27}
\]
At a split prime the rank-two Satake multisets concatenate;
at an inert prime one takes both square roots of each
rank-two parameter.  Formula (SX.1) shows that before the
last twist this attaches $R\chi^{-1}$; the $|\det|$ twist
attaches $R$.  At infinity, the two extensions of an embedding
of $L$ restrict to the two different real embeddings of $F_0$.
The infinitesimal character in (SX.27) is
\[
               \tfrac32+\{k+t,k,0,-t\}.
\]
Its four entries are distinct, giving regular algebraicity.

We verify every finite place in (SX.26).  At a place split in
$M/L$ the local parameters are the two rank-two principal-series
parameters restricted to the local extension.  Pullback by a
local norm preserves the real exponents in (SX.13), because
$q_w=q_v^f$ and
$|\nu(N_{M_w/L_v}\varpi_w)|=|\nu(\varpi_v)|^f$.
At a nonsplit place, the quadratic field $F_0$ is unramified
over the relevant local field: $\texttt{Two.unr}$ says its
quadratic character kills every original inertia subgroup.
Moreover this place is not above $2$ or $\ell$, since both
primes split in $F_0$.  The rank-two representation and the
rank-two local automorphic representation are therefore
unramified there, by SX1 and SX3.  Inducing an unramified
character through an unramified quadratic extension gives
the two unramified characters whose uniformizer values are
the two square roots.  Their absolute exponents are unchanged.
Thus in both cases the four local inducing characters have
the strict inequalities (SX.13), with one common center.
No quotient of two of them can equal $|\cdot|^{\pm1}$:
its real exponent lies strictly between $-1$ and $1$.
The principal-series induction is consequently irreducible
and has exactly these characters as its cuspidal support.
Finally the twist in (SX.27) subtracts $1$ from the center,
giving $s_4=s-1=-(k+3)/2$.
This proves the literal toral and strict-window clauses at
all finite places, including those above $\ell$.
\end{proof}

\subsection{Return to a finite coefficient ring and close the endpoint}

\begin{proposition}[SX7: exact finite coefficient return]
\label{claim:SX7}
The actual pair $(\delta R,(n_v))$ of SX5 descends to a
continuous pair $Q$ over the ring $O_{c_1}$ of one finite
extension $c_1\ge c$.  All its pair equations hold over that
ring.  Its actual reduction, rather than an abstractly
isomorphic representation, satisfies $\texttt{InV}$ and the
root-entry alternative in SX4.
\end{proposition}
\begin{proof}
Here is the matrix version of the finite-field capture argument
in RV4.  There are only countably many finite subextensions of
$\overline{\mathbb Q}_\ell/\mathbb Q_\ell$: in each degree
Krasner's theorem gives finitely many isomorphism classes, each
with finitely many embeddings.  Enumerate them inside an
increasing sequence $E_j$.  Each is complete and hence closed
in $\overline{\mathbb Q}_\ell$.
For the actual continuous homomorphism
$\rho=\delta R$, the sets
\[
           H_j=\{g\in G_F:\rho(g)\in\mathrm{GL}_4(E_j)\}
\]
are closed subgroups and cover $G_F$; inversion preserves $E_j$
because it is given by adjugate divided by the nonzero determinant.
The compact Hausdorff Baire theorem makes some $H_j$ have
interior, hence makes it open.  Choose its finitely many coset
representatives and adjoin all their matrix entries to $E_j$.
The homomorphism law now puts the entries of \emph{every}
$\rho(g)$ in this single finite extension.  Thus the argument
uses a closed-subgroup cover, not compactness of an arbitrary
algebraic-valued set.
Adjoin also the original coefficient field and the entries of the
finitely many matrices $n_v$ to this field.
This yields one finite extension.  SX4--SX5 show that all
matrix entries have absolute value at most one, so they lie
in its valuation ring.  The continuous topology there is
the adic topology, giving the literal continuity required
by $\texttt{SeedPairAlt}$.

Every pair equation is a polynomial identity with coefficients
in the original ring: the symplectic multiplier equation,
the killing of the frame inertia set, diagonalization at
$g_0$, the lower-unipotent equations for $n_v$, triangularity
on the original decomposition group, and triviality of its
diagonal on the original $\operatorname{muOf}(D_v)$ set.  Injectivity
of $O_{c_1}\hookrightarrow\overline{\mathbb Q}_\ell$
therefore descends them all.
Reduction also commutes with the symplectic adjoint and
conjugation.  An entry reduces to zero exactly when its
absolute value is less than one; a unit entry remains
nonzero.  Thus (SX.14), (SX.23)--(SX.24), and their unit
$(2,0)$ entries give literal $\texttt{InV}$.
The entries $(1,0)$, or $(1,2)$ and $(3,0)$, give respectively
$\operatorname{Ent}(-w_\alpha)$, or
$\operatorname{Ent}(w_\beta)$ and
$\operatorname{Ent}(-w_\vartheta)$.
The diagonal character congruences and the local ordered
character congruences give the original residual pair record.
All these assertions concern the descended matrices themselves.
\end{proof}

\begin{theorem}[SX8: the literal \texttt{SeedIIAlt} supplier]
\label{claim:SX8}
For every original $\Omega$, prime $\ell$, coefficient datum
$c$, raw datum $D$, and tower $T$ satisfying (SX.2), and for
every fixed $\iota:\overline{\mathbb Q}_\ell\simeq\mathbb C$,
there exist a finite coefficient extension $c_1\ge c$ and an
actual pair $Q$ over $O_{c_1}$ such that
\[
 \begin{gathered}
 \texttt{SeedPairAlt}(\iota,c\le c_1,Q),\qquad
 \texttt{InV}(\bar Q),\\
 \operatorname{Ent}(-w_\alpha,\bar Q.r)\quad\text{or}\quad
 \bigl(\operatorname{Ent}(w_\beta,\bar Q.r) 
       \text{and}\ \operatorname{Ent}(-w_\vartheta,\bar Q.r)\bigr).
 \end{gathered}                                                \tag{SX.28}
\]
The only retained leaves of this construction are the precisely
stated classical theorems above and the already reconstructed
coefficient-capture theorem.  In particular no mixed-seed,
ordinary-weight-change, or automorphic-witness predicate remains
as an unproved bespoke supplier.
\end{theorem}
\begin{proof}
Take the pair of SX7.  It is a pair for the unchanged frame
$T.\mathrm{fr}(\varnothing)$, is continuous, and has its
literal residual pair and (SX.28), by SX4--SX7.
Twisting by a scalar character does not change invariant
subspaces, so SX6 gives irreducibility of $\delta R$ on
the original $T.G'=G_L\cap\ker\chi$.
Use the weights $h_i=2m_i+b_0$ of (SX.20) and put
\[
          \kappa=\delta^{-1}|_{G_L},\qquad a=-b_0.
                                                            \tag{SX.29}
\]
Then $h_i-h_{i+1}=2(m_i-m_{i+1})$ with each difference
$m_i-m_{i+1}\ge2$, and $a+h_0=2m_0$ is even.
The character $\kappa$ is continuous.  For the same $N>0$
from (SX.21),
\[
       (\kappa(s)^2\chi(s)^{-a})^N=1
\]
for every inertia element at every valuation above $\ell$;
in fact it holds for every $s\in G_L$.
Twisting $(\delta R)|_{G_L}$ by $\kappa$ gives exactly
$R|_{G_L}$, so SX6 supplies the actual $\Pi$, the attachment,
regular algebraicity, and the required $\texttt{AllWin}$.
These are all clauses of $T.\texttt{IsSeedAlt}$.

The order of choices is the literal one: $D,T$ and all their
data are given first, then $\iota$ is given, and only then are
the Hilbert form, $r$, $R$, $\delta$, $\Pi$, and $c_1,Q$
chosen.  The argument applies to every such given tower.
It does not choose a new tower after seeing the seed, or
require a new complex embedding at a later stage.
\end{proof}

\subsection{Primary reference locations and scope of the closure}

The classical references used above are the following.  The accompanying
machine-readable registry records their individual applications and
the exact pinned Lean definitions replaced or matched by this chapter.
The papers themselves are research references and are not redistributed.

\begin{itemize}
\item S.~Dasgupta, H.~Darmon and R.~Pollack,
\emph{Hilbert modular forms and the Gross--Stark conjecture},
Annals of Mathematics 174 (2011), 439--484,
\S2.2, Proposition 2.1, especially formula (50), pp.~458--459;
\S2.3, p.~466, for the usual cuspidal/Eisenstein decomposition.
\href{https://annals.math.princeton.edu/wp-content/uploads/annals-v174-n1-p12-p.pdf}{Primary article (PDF)}.
The later Gross--Stark hypotheses are not used in the general
Eisenstein-series proposition.

\item C.~P.~Mok and F.~Tan,
\emph{Overconvergent Families of Siegel--Hilbert Modular Forms},
Canadian Journal of Mathematics 67 (2015), 893--922,
Lemma 2.10 and Proposition 2.11, pp.~905--906.
\url{https://doi.org/10.4153/CJM-2014-017-9}.
The genus-one determinant Hasse invariant and Serre lifting are
the only results used; no eigenvariety or overconvergent
automorphy assertion is imported.

\item H.~Hida, \emph{Control of nearly ordinary Hecke algebras},
Luminy lecture, 18 July 2007, pp.~3--8,
Theorem 1.2, Lemmas 2.1--2.2, Theorem 2.4,
Corollary 2.5 and Theorem 3.1 with Remark 3.1.
\url{https://www.math.ucla.edu/~hida/Lum07.pdf}.
The central character is fixed, $F_0$ is real quadratic,
$\ell$ splits, and only characteristic-zero classical
specializations are needed.

\item J.~Arthur and L.~Clozel,
\emph{Simple Algebras, Base Change, and the Advanced Theory
of the Trace Formula}, Annals of Mathematics Studies 120 (1989),
Chapter III, Theorems 4.2, 5.1 and 6.2,
pp.~202--203, 212--213 and 215--217.
\url{https://www.claymath.org/library/cw/arthur/pdf/30.pdf}.
For strong local compatibility also use G.~Henniart,
\emph{Induction automorphe globale pour les corps de nombres},
Bulletin de la Soci\'et\'e Math\'ematique de France 140 (2012),
1--17, Theorems 4--5 and \S3.7, pp.~6--7 and 16:
\url{https://www.numdam.org/article/BSMF_2012__140_1_1_0.pdf}.
Their unitary formulations apply after removing the displayed
common norm twist; base change, induction and restoring that
twist commute by their local character formulas.

\item T.~Barnet-Lamb, T.~Gee, D.~Geraghty and R.~Taylor,
\emph{Potential automorphy and change of weight},
Annals of Mathematics 179 (2014), 501--609,
\S2.1 and Theorem 2.1.1(1)--(2), pp.~535--538.
\href{https://annals.math.princeton.edu/wp-content/uploads/annals-v179-n2-p03-p.pdf}{Primary article (PDF)}.
Only the rank-two polarized CM Galois attachment and
away-from-$\ell$ compatibility are used.  The actual Hilbert
form has already been constructed before applying this theorem.

\item W.~Casselman, \emph{On some results of Atkin and Lehner},
Mathematische Annalen 201 (1973), 301--314,
Theorem 1 and its conductor computation, pp.~302--306.
\url{https://lesesvre.perso.math.cnrs.fr/newforms-references/casselman.pdf}.
The use is the exact equivalence between a prescribed
newvector level and the local conductor.  The principal-series
classification and its determinant/conductor normalization are
also in H.~Jacquet and R.~P.~Langlands,
\emph{Automorphic forms on GL(2)}, Chapter I, \S3:
\url{https://publications.ias.edu/sites/default/files/automorphic-forms-on-gl2_rpl_9.pdf}.

\item I.~N.~Bernstein and A.~V.~Zelevinsky,
\emph{Induced representations of reductive $p$-adic groups. I},
Annales scientifiques de l'\'Ecole Normale Sup\'erieure 10
(1977), 441--472, Theorem 4.2, p.~454.
\url{https://www.math.tau.ac.il/~bernstei/Publication_list/publication_texts/BZ-induced-1.pdf}.
We use its sufficient irreducibility criterion only for
four one-dimensional inducing characters.

\item J.~S.~Milne, \emph{Class Field Theory}, version 4.03,
Chapters I and VIII, for local and global reciprocity;
\emph{Algebraic Number Theory}, version 3.08,
the Dirichlet unit theorem and Theorem 8.31 (Chebotarev):
\url{https://www.jmilne.org/math/CourseNotes/CFT.pdf},
\url{https://www.jmilne.org/math/CourseNotes/ANT.pdf}.
The resulting elementary character and irreducibility arguments
are written out in SX1--SX6.
\end{itemize}

This closes the literal $\texttt{SeedIIAlt}$ supplier with its
original $\texttt{Ctx2}$, $\texttt{Ok}$, $\texttt{Two}$ and
$\texttt{Deep}$ hypotheses.  It does not claim the wider
$\texttt{HypothesesH'}$ or $\texttt{RankN}$ packages, general
rank-four ordinary support, $\texttt{SeedIVAlt}$, image density,
or the final counting theorems.  Those are separate dependency
edges, and the fixed $\iota$ in this endpoint remains the one
chosen before the seed.
